\section{Pengujian dan Pembahasan}
\label{sec:pengujian}


Pada bagian ini mahasiswa melakukan pengujian terhadap rancangan yang diajukan di \ref{chap:Detail_Implementasi}. Mahasiswa harus menuliskan hal-hal diantaranya: 
\begin{itemize}
    \item Tujuan pengujian,
    \item Hal atau bagian apa yang diuji,
    \item Skenario pengujian,
    \item Analisis dari hasil pengujian. Analisis harus dikaitkan dengan ilmu-ilmu dasar (matematika, fisika, dan \textit{engineering science}) kenapa suatu fenomena terjadi. Analisis juga merupakan cara untuk memastikan rancangan yang diajukan di \ref{chap:Detail_Implementasi} memenuhi spesifikasi atau kriteria sukses yang disebutkan,
    \item Hal-hal penting lainnya.
\end{itemize}

\subsection{Skenario Pengujian 1 dan Analisis}
\label{subsec:Skenario1}
Sesuaikan jumlah skenario sesuai kebutuhan pengujian anda
\subsection{Skenario Pengujian 2 dan Analisis}
\label{subsec:Skenario2}
Sesuaikan jumlah skenario sesuai kebutuhan pengujian anda

\section{\textit{Improvement}}
\label{sec:improvement}
Di dalam Bab ini, Anda harus secara rinci menjelaskan kenapa perlu tambahan dan/atau revisi rancangan secara logis dan kronologis sehingga tim penilai akan melihat bahwa tambahan ini adalah sesuatu yang perlu dan wajar untuk dilakukan. 


